\chapter{Étude de l’existant}
\section{Objectif et méthode}
L’étude de l’existant vise à distinguer les fonctionnalités déjà présentes sur le marché tunisien des éléments différenciants envisagés pour Souk El Hayawan. Elle repose sur la consultation des pages publiques des plateformes recensées dans le document de cadrage le \textbf{19 octobre 2026}. Les constatations sont limitées aux fonctions présentées ou observables ; elles ne présument pas de l’architecture interne des sites.

\section{Plateformes étudiées}
\subsection{Tayara --- annonces généralistes}
Tayara réunit des annonces de nombreux domaines, dont une rubrique consacrée aux animaux. Elle offre un parcours généraliste de publication et de recherche d’offres. Cette polyvalence est un atout, mais elle ne correspond pas à une spécialisation exclusive sur les fiches propres aux animaux d’élevage et de compagnie.\par
\source{\url{https://www.tayara.tn/}}

\subsection{Sawt El Falleh --- marché agricole}
Sawt El Falleh s’adresse au secteur agricole tunisien. Les fonctionnalités décrites dans le cadrage comprennent des annonces de vente et d’achat, des catégories, des filtres régionaux et de prix, ainsi que le contact direct. Il s’agit d’un point de comparaison important : \textbf{rechercher un animal par prix et par région n’est pas en soi une innovation}.\par
\source{\url{https://www.sawtelfalleh.tn/}}

\subsection{TunisieVet --- annonces et orientation vétérinaire}
TunisieVet propose un espace d’annonces d’animaux et des contenus associés aux services vétérinaires. Son formulaire de vente prévoit des informations comme l’espèce, la race, l’âge, le poids, les photos, le prix et la localisation. La structuration des données animales existe donc également dans cette offre.\par
\source{\url{https://tunisievet.com/fr/vente-animaux}}

\subsection{Animalia --- animaux de compagnie}
Animalia est orientée vers les animaux de compagnie et l’adoption responsable. Ses pages de présentation mettent en avant les annonces et des engagements liés à la vérification. Cette orientation est directement pertinente pour la partie « compagnie » de notre futur produit.\par
\source{\url{https://www.animalia.tn/}}

\section{Comparaison qualitative}
\begin{table}[htbp]
\centering\small
\caption{Positionnement des plateformes étudiées}
\label{tab:comparatif}
\begin{tabularx}{\textwidth}{@{}P{31mm}Y Y@{}}
\toprule
\textbf{Plateforme} & \textbf{Orientation relevée} & \textbf{Apport pour notre conception} \\
\midrule
Tayara & Petites annonces généralistes & Comprendre un parcours de consultation multi-catégories.\\
Sawt El Falleh & Échanges agricoles ; filtres catégorie, région et prix & Consolider des fiches animales plus adaptées aux espèces.\\
TunisieVet & Annonces animales et informations vétérinaires & S’inspirer de données structurées par animal.\\
Animalia & Annonces et adoption d’animaux de compagnie & Étudier un parcours d’adoption et de modération.\\
\bottomrule
\end{tabularx}
\end{table}

\begin{warningbox}
\textbf{Limite de l’étude :} l’absence d’une fonctionnalité sur une page visitée ne prouve pas son absence de la plateforme. Il serait incorrect d’affirmer qu’aucun concurrent ne propose la recherche par prix, les races ou les filtres géographiques.
\end{warningbox}

\section{Opportunité identifiée}
Le positionnement envisagé repose sur la \textbf{cohérence d’un parcours multi-espèces} plutôt que sur une fonctionnalité isolée : annonces adaptées à la nature de l’offre, données comparables, filtres combinables, prise en charge des régions tunisiennes, et espace spécifique d’entraide. Une étude utilisateur complémentaire reste nécessaire pour confirmer que cet ensemble apporte une valeur réellement perçue.

\section{Conclusion de l’étude}
L’étude montre un marché déjà occupé par des services généralistes ou spécialisés. Souk El Hayawan devra donc être évalué sur la pertinence de ses parcours, la qualité de ses données, la simplicité de l’expérience et sa modération, plutôt que sur une prétention de première plateforme du secteur.
